\section{Structure of the proof}
\label{sec:proof-structure}

The complement quotient explains the endpoint alternatives in the
three-prime classification. Its low-root bound localizes a possible
integer spectrum to two endpoint surfaces, whose geometric bounds lead
to different spectral obstructions. We summarize this mechanism before
giving the component proofs.

\subsection{Localization and the sharp threshold}

Let $C$ be the complement Laplacian quotient on the six proper nonempty
supports, of weights $a,b,c,ab,ac,bc$, and put $h_C(x)=\det(xI-C)/x$.
The off-diagonal entry from $S$ to $T$ is $-w_T$ when $S\cap T$ is empty,
and zero otherwise; the diagonal is the sum of these neighbor weights.
The weighted symmetric quotient is positive semidefinite with a simple
zero eigenvalue, since the disjointness graph on these supports is connected.

For positive integer exponents, every nonzero complement quotient root
must be an integer if $G_n$ is Laplacian integral: reflection across the
complement and the universal-vertex lift transfer every noninteger root
to a noninteger graph eigenvalue. Theorem~\ref{thm:uniform-low-root}
gives a positive quotient root below three for every ordered positive
real triple. Hence the smallest positive root also lies in $(0,3)$,
and an integral graph spectrum requires
\[
 h_C(1)=0\qquad\hbox{or}\qquad h_C(2)=0.
\]
These are necessary conditions; the following arguments exclude integer
spectra on each surface by a further root or an exponent-root bracket.

\begin{proposition}
\label{prop:localization-sharpness}
The uniform threshold three is optimal even for integer exponents.
For $a=b=c=m\geq1$, the smallest positive quotient root satisfies
\[
 \mu_-=
 \frac{m^2+4m-m\sqrt{(m+4)^2-12}}{2}
 \in\left(\frac{3m}{m+4},3\right),
 \qquad \lim_{m\to\infty}\mu_-=3.
\]
\end{proposition}
\begin{proof}
Direct substitution in the quotient gives
\[
 h_{m,m,m}(x)=(x-m^2-m)
 \bigl[x^2-(m^2+4m)x+3m^2\bigr]^2.
\]
For the quadratic $q_m$, we have $q_m(0)>0$ and $q_m(m)=-m^3<0$.
Its smaller root therefore lies in $(0,m)$ and is below the other
three positive roots, counted with multiplicity. Also
$q_m(3)=9-12m<0$, so $\mu_-<3$. Rationalization gives
\[
 \mu_-=
 \frac{6m}{m+4+\sqrt{(m+4)^2-12}}
 >\frac{3m}{m+4}.
\]
Dividing the denominator by $m$ proves the limit, including along
positive integers. Thus no smaller constant is a uniform upper bound.
\end{proof}

\subsection{The endpoint surfaces and their scales}

After repeated exponents and minima at most seven have been settled,
the remaining classification domain is $8\leq a<b<c$. In this domain,
Lemma~\ref{lem:endpoint-two-rectangle} shows that
\[
 h_C(2)=0\quad\Longrightarrow\quad
 b<2a-2,\qquad c<\frac{9a}{4}-8.
\]
Both larger exponents are bounded multiples of the minimum. The
rectangle is finite at each fixed integer $a$, but $a$ is unbounded.
Theorem~\ref{thm:endpoint-two-completion} supplies a root in $(4,5)$
for $a\geq40$; the complete $8\,658$-triple certificate handles
$8\leq a\leq39$. The real tail argument and the finite base together
close this alternative.

Endpoint one has a different necessary scale. For real
$2\leq a\leq b\leq c$, Theorem~\ref{thm:endpoint-one-spectrum} gives
\[
 h_C(1)=0\quad\Longrightarrow\quad
 b+c\geq2a^2-2a+1,\qquad c\geq a^2-a+\tfrac12.
\]
The second inequality follows from $b\leq c$. This quadratic lower
bound supplies the hypothesis of the largest-root interval in
Theorem~\ref{thm:largest-root-unit-interval} whenever $b-a\geq3$.
For gaps one and two, Lemma~\ref{lem:endpoint-one-small-gaps} instead
places the only possible exponent-root above $b$ between consecutive
integers. This route settles the spectrum without enumerating integer
points on the endpoint-one surface.

\subsection{The fixed-gap endpoint-one expansion}

The small-gap anchors also have an asymptotic interpretation. The
following result concerns the exponent-root $c$, with the gap fixed;
it is supplementary to the exact brackets used in the classification.

\begin{proposition}
\label{prop:fixed-gap-expansion}
Fix a real $d\geq1$ and put $b=a+d$. For all sufficiently large real
$a$, the equation $h_C(1)=0$, regarded as a cubic in $c$, has exactly
one root $c^*$ above $b$. As $a\to\infty$ with $d$ fixed,
\begin{align}
 c^*={}&2a^2+(2d-3)a-\frac{(d+2)(2d-1)}2
 \notag\\
 &+\frac{d^2}{2a}
 +\frac{d^2(d+1)(2d-3)}{4a^2}+O_d(a^{-3}).
 \label{eq:fixed-gap-expansion}
\end{align}
The remainder constant and the required minimum may depend on $d$.
\end{proposition}
\begin{proof}
Write $Q_{a,d}(c)=h_C(1)$ at $b=a+d$. The cubic coefficients are
displayed in Lemma~\ref{lem:endpoint-one-small-gaps}; their identity
holds for arbitrary $d$. Set $t=1/a$, $y=t^2c$ and
\[
 F_d(t,y)=t^8Q_{1/t,d}(y/t^2).
\]
This is a polynomial in $t,y$, with
\[
 F_d(0,y)=2y^2(y-2),\qquad
 \partial_yF_d(0,2)=8.
\]
The analytic implicit-function theorem gives a unique real analytic
branch $y_d(t)$ near $t=0$ with $y_d(0)=2$. Substitution and comparison
of successive powers of $t$ give
\begin{align*}
 y_d(t)={}&2+(2d-3)t-\frac{(d+2)(2d-1)}2t^2
 +\frac{d^2}{2}t^3\\
 &+\frac{d^2(d+1)(2d-3)}4t^4+O_d(t^5).
\end{align*}
Multiplication by $t^{-2}$ proves~\eqref{eq:fixed-gap-expansion}.
For fixed $d$ and sufficiently large $a$, the cubic has positive
leading coefficient and $Q_{a,d}(0)>0$, whereas
\[
 Q_{a,d}(a+d)=-4a^6+O_d(a^5)<0.
\]
It therefore has a root in each of $(-\infty,0)$ and $(0,b)$.
The analytic branch has $c^*\sim2a^2>b$ and supplies the third root.
These distinct roots exhaust the cubic degree, proving uniqueness
above $b$.
\end{proof}

For $d=1,2$, the polynomial parts of this expansion are respectively
$2a^2-a-2$ and $2a^2+a-6$, followed by a positive term of order
$a^{-1}$. The exact unit brackets in
Lemma~\ref{lem:endpoint-one-small-gaps} hold for all real $a\geq2$.
The expansion here is for fixed $d$; it supplies no uniform error bound
when $d$ grows with $a$. The largest-root argument handles every
$d\geq3$ in the classification using the quadratic lower bound on $c$.
